\cotefile{12}
\scanpage{51}
$X$ \tikz[baseline=(a.base)]\node[ellipse,draw,inner sep=2pt](a){lisse}; \deleted{\unsure{géom.|équidim.} \unsure{conn.}} de dim $n$ partout

\[
\left\{\begin{aligned}
H^i_{!}(X,\mathbb Z(0))&\in\operatorname{Ob} M_i^+(K)\,\deleted{\unsure{\operatorname{Ob}}}\ \inserted{\text{et }\in\operatorname{Ob}M_n^+(K)}\quad\text{tt }i\qquad\text{mieux :}\\
H^{n+k}_{!}(X,\mathbb Z(0))&\in\operatorname{Ob}\big[\unsure{\zeta|\xi}^{k}M_{n-k}^{+}(K)\big]\deleted{+\zeta^{k+1}M^{+}_{n-k-2}}
\end{aligned}\right.
\]

\[
H^i(X,\mathbb Z(0))\in\operatorname{Ob}
\underbrace{D_i(M^+(K)),\ \text{et }\in\operatorname{Ob}D_n(M^+(K))}_{\text{et }=\text{ mieux}}\qquad\text{tt }i
\]
\[
\begin{aligned}
&\deleted{H^i(X,\mathbb Z(0))\in}\\
&H^i(X,\mathbb Z(0))\in\operatorname{Ob}
\big[\deleted{\unsure{\zeta}}M_i^++\zeta M_{i-1}^++\zeta^2 M^+_{i-2}+\cdots\big]\qquad\text{si }i\leq n\\
&\hspace{6em}\deleted{M_n^++\zeta M_n^++\zeta^2M_{n-2}^++\cdots}\\
&\hspace{3em}\in\operatorname{Ob}\zeta^{i-n}
\big[M_{2n-i}^++\zeta M_{2n-i-1}^++\zeta^{2}M_{2n-i-2}^+\ -\big]\qquad\text{si }i\geq n
\end{aligned}
\]

\[
\operatorname{cl} H^i(X,\mathbb Z(0))\in\sum_{1\leq\alpha\leq\operatorname{Inf}(2i,2n)}\operatorname{Gr}_i(M^+(K))
\qquad\text{plus gén.}
\]
\[
\operatorname{cl} H^i(X,M)\in\sum_{p+1\leq\alpha\leq p+\operatorname{Inf}(2i,2n)}\operatorname{Gr}_i\qquad\text{si } M\ \deleted{\unsure{pur|pure}}\ \text{de poids }p
\]
% archive sheet number: 50

\scanpage{52}
% One long diagonal stroke runs through the entire written area.
\deleted{(i.e. l'Albanese à isogénie près) d'une surface $X$ (sur \unsure{$k$|$K$} alg. clos) soit connu quand on \unsure{connait|connaît} $X$ \unsure{ds|dans} $L(K)$ seulement \ldots
\[
\deleted{L_i(K)/L_{i-1}(K)\longrightarrow M_i^+(K)/M_{i-1}^{\unsure{+}}(K)}
\]
Mettons que l'hom
\[
L_i(K)/L_{i-1}(K)\longrightarrow D_i M^+(K)/D_{i-1}M^+(K)
\]
prend ses valeurs dans l'image de $M_i^+(K)$ modulo $D_{i-1}M^+(K)$, \deleted{\unsure{da|donc}} \inserted{i.e.} dans la partie primitive. On peut se demander \deleted{\unsure{si}} \inserted{si}
\[
\boxed{L_i(K)/L_{i-1}(K)\longrightarrow\deleted{\unsure{\sigma_i(K)}}\ \mathbb Z^{\sigma_i(K)}.}
\]
est injectif, ou bijectif, du moins mod torsion\ldots.}
% archive sheet number: 51

\scanpage{53}
\mspage{\textcircled{13}}
\deleted{\uline{12}}\quad\uline{Caractérisation galoisienne des filtrations}

\margin{\unsure{13|B}}
On peut définir le sous-groupe (ou \unsure{pos-anneau}) $\mathrm{M}^+_\ell(X)$ du $\unsure{K|\mathbb K}$ ($K$ des \deleted{\unsure{$\mathbb Q_\ell$-}\illegible{1}} $\mathbb Q_\ell$-faisceaux constants tordus sur $X$) \uline{provenant} de motifs, on a donc un hom surjectif
\[
M^+(X)\longrightarrow M^+_\ell(X)
\]
\inserted{Si $X$ est de type fini sur $\mathbb Z$,} Cet hom est en fait \uline{bijectif}, car il résulte des conjectures de Tate, plus précisément, que si $E$ et $F$ sont des motifs constants tordus, \uline{simples} \inserted{\uline{non isom.}} sur $X$, alors $T_\ell(E)$ et $T_\ell(F)$ \deleted{\unsure{sont}} n'ont pas de \unsure{facteurs|facteur} simples isomorphes communs.

Par suite, \deleted{\unsure{donc}} la \deleted{\unsure{fil|filt.}} \uline{graduation} de $M^+(\unsure{x|X})$ (via une partition de $\Sigma^+(x)$) et la filtration de $M^+(x)$ (via une filtration de $\Sigma^+(x)$) peuvent se décrire en termes galoisiens. On trouve \inserted{pour $E\in\operatorname{Ob} M(X)$} :

\deleted{$\operatorname{cl}(E)$ de degré $i$ \unsure{s} $\in$}
% slanted
\margin{``Conjectures de Weil\textsuperscript{\unsure{\#|de}} Riemann''}
\[
\left\{\begin{aligned}
\inserted{\operatorname{cl}(E)\in{}}M_i(X)&\Longleftrightarrow\ \text{pour tt pt fini }x\text{ de }X,\\
&\text{les valeurs propres de Frobenius dans }T_\ell(E)(x)\text{ sont}\\
&\text{de val. abs. }N(x)^{i/2}
\end{aligned}\right.
\]
% archive sheet number: 52

\scanpage{54}
D'autre part, si $E\in\operatorname{Ob}\unsure{\mathcal M|M}^{\deleted{\unsure{+|\ast}}}(X)$, \deleted{\unsure{p. ex. simple|simple}}, on trouve
% slanted
\margin{$\nearrow$ ``Conjectures de Tate'' $\searrow$}
\[
\left\{\begin{aligned}
\operatorname{cl}(E)\in M^+(X)\quad\Longleftrightarrow\quad
&\forall\,x\ \text{pt fini de }X,\ \text{le polynôme}\\
&\text{caractéristique }\deleted{\text{de Frobenius}}\\
&\operatorname{Frob}_x^{-1}\ \text{dans }T_\ell(E)(x)\text{ est}\\
&\text{à coefficients }\text{\uline{entiers}}\\
&[\,\text{i.e. les valeurs propres de}\\
&\text{Frobenius sont des}\\
&\text{\uline{entiers} algébriques}\,]
\end{aligned}\right.
\]

Par ailleurs, pour \deleted{\unsure{un élément, pour|décrit, pour} $\unsure{\zeta^\alpha}\in M_\alpha(X)$} $E$ ``homogène'' de degré $\alpha$, pour exprimer si $E\in\operatorname{Ob}D_i(M(X))$, où $i\geq 0$ \deleted{on peut, \unsure{par}}
\inserted{\deleted{$\operatorname{cl}(E)\in\unsure{\Sigma^i|\sum_i}$} \deleted{$\unsure{\zeta^\alpha}\in D_\alpha M(X)$}
$E\in\operatorname{Ob}D_i\mathcal M(X)$ \unsure{\}} i.e. $\operatorname{cl}(E)\in D_i M(X)\Longleftrightarrow$ Si} on est $\Downarrow$

$[\,1^\circ)$ $i=0$. Alors
\[
\begin{aligned}
E\in\operatorname{Ob}D_0(M(X))\Longleftrightarrow\quad &E\text{ est de degré \uline{pair} }\alpha=2\beta\\
&\text{et }E\unsure{\zeta|\xi}^{-\beta}=F\quad(\text{qui est de}\\
&\text{degré }0)\text{ est }\in\operatorname{Ob}M^+(X)\\
&\text{i.e. si }\rho_{\unsure{\nu|i}}(x)\text{ sont les}\\
&\text{valeurs propres de }\operatorname{Frob}_x^{-1}\text{ dans}\\
&T_\ell(M(x)),\ \text{alors les }\rho_\nu(x)/N(x)^\beta\\
&\text{sont des entiers algébriques}\,]
\end{aligned}
\]
\margin{N.B. Ce sont bien des racines de l'unité [leur produit étant $\pm1$ !]}

$2^\circ)$ $i\geq 1$. Alors il existe une plus petite puissance $\zeta^\beta$ de $\zeta$ telle que $E\zeta^{-\beta}\deleted{\unsure{\in\operatorname{Ob} M}}\inserted{{}=E(+\beta)}$ soit de degré $i-1$ ou $i$ \unsure{[|(}suivant la parité de $\alpha\,]$. On veut que \deleted{\unsure{cet}} $E(\unsure{+\beta|\beta})$ soit $\in\operatorname{Ob} M^+(X)$, i.e. les $\rho_\nu(x)/N(x)^\beta$ doivent être des entiers alg.
% archive sheet number: 53

\scanpage{55}
Le cas $1^\circ)$ \deleted{\unsure{doit|dit}} signifie aussi que $T_\ell(E)(\beta)$ est tel que $\pi_1(X)$ y opère à travers un \uline{quotient} \uline{fini} [Mais cet énoncé ne se généralise guère via filtration de $\pi_1$, comme on voit sur $X=\operatorname{Spec}$ corps fini].

Supposons $E$ \inserted{$\in\operatorname{Ob}M^+$} \uline{semi-simple}, donc $T_\ell(E)$ semi-simple. Soit $E_i$ la composante de $E$ de type dimensionnel $i$, \deleted{\unsure{qu}} comment caractériser $T_\ell(E_i)\subset T_\ell(E)$ du pt de vue Galoisien ?
\inserted{Lorsque $X=\operatorname{Spec}$ corps fini, c'est purement \unsure{arithmétique|combinatoire} en termes de Frobenius.}
\deleted{On \unsure{pt|peut} supposer $E$ homogène de degré $i$ [la décomposition de $T_\ell(E)$ en composantes \deleted{\unsure{de degr.}} homogènes est canoniquement déterminée]. On est ramené de proche en proche \deleted{[tordant par \unsure{$\mathbb Z$} $\mathbb Z(1)$ et} [\unsure{dévissant|diminuant} par le type dimensionnel précédent] à trouver le morceau de type \inserted{dimensionnel} $i$, qd on a de type dimensionnel \deleted{$0$ ou de type dimensionnel $1$ (alors de type dimensionnel $0$ (qui sont en \unsure{qui sont nuls}}}
Le cas général s'en déduit par image inverse, si on connaît \uline{tous} les Frobenius $\in\pi_1$ \unsure{au}
% archive sheet number: 54

\scanpage{56}
14) \uline{Invariants de Galois et \unsure{thmes|thms} de commutation}

$X$ propre lisse sur $K$ de type fini

$H^{2i}(\bar X,\deleted{\mathbb Z}\mathbb Q_\ell(i))$ est de poids $0$, donc peut avoir des invariants non nuls sous Galois.

La conjecture de Tate s'énonce en disant que
\[
H^{2i}(\bar X,\deleted{\mathbb Z}\mathbb Q_\ell(i))^\pi
\xleftarrow{\sim}\unsure{\mathscr A|A|\mathcal A}^{i}(X/K)\otimes_{\mathbb Q}\mathbb Q_\ell
\]

On peut \unsure{le|la} préciser ainsi :

L'anneau des correspondances alg $/K$ \deleted{\unsure{en}}, \inserted{opérant en dim $i$} \deleted{\unsure{$C^i(X)$}} $\unsure{\mathscr U|\mathscr W|\mathcal U}^i(X/K)$ est un anneau semi-simple, et [via Hodge] et on veut que
\[
\begin{aligned}
&\mathscr U^i(X/K)\otimes_{\mathbb Q}\mathbb Q_\ell\ \text{soit le commutant de }\pi=\operatorname{Gal}(\bar K/K)\\
&\hspace{9em}\text{dans }\operatorname{End}_{\mathbb Q_\ell}(H^i(X,\mathbb Q_\ell(0))),
\end{aligned}
\]
\uline{et} l'alg. enveloppante de $\pi$ ds $\operatorname{End}_{\mathbb Q_\ell}(H^i)$ doit être le commutant de $\mathscr U^i(X/K)$ [donc \uline{l'action de $\pi$ est $1/2$ simple}]

Le deuxième énoncé $+$ injectivité de $\mathscr U^i(X/K)\otimes_{\mathbb Q}\mathbb Q_\ell \longrightarrow\operatorname{End}_{\mathbb Q_\ell}(H^i(X,\mathbb Q_\ell(0)))$ entraîne tout.
% archive sheet number: 55

\scanpage{57}
\uline{N.B.} a) Ceci ne caractérise pas les classes algébriques elles-$=$s, mais seulement \unsure{le|les} $\mathbb Q_\ell$-\unsure{espace|espaces} \unsure{engendré|engendrés}.

b) De $=$, on ne peut caractériser ainsi l'image de $\pi$, mais seulement l'algèbre enveloppante\ldots.

Passant à la limite sur les extensions finies de $K$, on trouve
\[
\left\{\begin{aligned}
&H^{2i}(\bar X,\mathbb Q_\ell(i))^{\unsure{\mathfrak g|g}_\ell(\bar K/K)}
\xleftarrow{\sim}\unsure{\mathscr A|A|\mathcal A}^{i}(\bar X)\otimes_{\mathbb Q}\mathbb Q_\ell\\
&\unsure{\mathscr U|\mathscr W|\mathcal U}^i(X)\otimes_{\mathbb Q}\mathbb Q_\ell\quad\text{et l'algèbre enveloppante de}\\
&\hspace{8em}\mathfrak g_\ell(\bar K/K)\ \text{ds }\operatorname{End}_{\mathbb Q_\ell}(H^i(\bar X,\mathbb Q_\ell(0))\\
&\hspace{8em}\text{sont commutants l'un}\\
&\hspace{8em}\text{de l'autre.}
\end{aligned}\right.
\]

\uline{Question} Comment caractériser l'image de $\mathfrak g_\ell(\bar K/K)$ dans $\operatorname{End}_{\mathbb Q_\ell}(H^i(\bar X,\mathbb Q_\ell(0))$ ??

\unsure{Cela|Cel} doit être une sous Algèbre de Lie \uline{réductive}, on pt la considérer comme plongée ds $H^{2n-i}(\bar X)\otimes H^i(\bar X)(n)$.
% archive sheet number: 56

\scanpage{58}
\deleted{dimensionnel $\leq n-2p$, donc que les $\operatorname{gr}^{p}_{\unsure{F|f}}(H^n(X,\mathbb Z(0)))$ sont \deleted{\unsure{du\ que\ le\ poids|bipurs}} de type \uline{dimensionnel} \uline{purement} \uline{$n-2p$}. Or de façon précise, on sait qu'il est de \uline{poids} $n$, d'autre part le terme
\[
E_1^{pq}=\coprod_{x\in X(p)} H^{n-2p}(x,\mathbb Z)(-p)
\]
\deleted{\unsure{(pour}} (\unsure{ici|où} $q=n-p$) \unsure{chacun|donne} des $M^+(X)$ \unsure{en} \unsure{sont|sot} dans
\[
\unsure{\zeta|\xi}^{p}\big[M_{n-2p}^++\zeta M_{n-2p+1}^++\cdots\big]
\]

% margin note written sideways
\diagram{\begin{tikzcd}[ampersand replacement=\&,column sep=large,row sep=small]
H^j_Y(X)\arrow[r,dash] \& H^j(X)\arrow[r,dash] \& H^j(U)\arrow[r,dash] \& H^{j+1}_Y(X)\\
H^i_!(Y) \& H^i_!(X)\arrow[l] \& H^i_!(U)\arrow[l] \& H^{\deleted{\unsure{i}}i-1}_!(Y)\arrow[l]
\end{tikzcd}}
\[
\begin{aligned}
&i+j=2n\qquad j\leq n\\
&H^i_!(Y)\in\zeta^{i-m} M_{2m-i}^{+}\qquad i\geq n>\dim Y=m\\
&\hspace{13em}\deleted{\unsure{i+j}}\\
&\hspace{3em}\cap\\
&\zeta^{i-m}M_{2n-i-2}^{+}=\zeta^{i-\deleted{\unsure{m}}m}M_{j-2}^{+}\\
&H_!^{i-1}(Y)\in\zeta^{i-m-1}M_{2m-(i-1)-2}^{+}=\zeta^{i-m-1}M_{j-1}^{+}
\end{aligned}
\]
si $i$}
% archive sheet number: 57

\scanpage{59}
Si $M$ est un motif sur $K$, $1/2$ simple pour simplifier, alors on a
\[
\unsure{\mathfrak g|g}_\ell(\overset{\unsure{\ell|s}}{\bar K}/K)\longrightarrow\operatorname{End}(M_l)\simeq\check M_l\otimes M_l
\]
On pense que l'image provient d'un motif $\mathfrak g_{\deleted{(M)}M}\inserted{{}\subset\check M\otimes M}$ indépendant de $l$, [qui est un motif en \unsure{Algèbre|algèbre|algèbres} de Lie].

Pour $M$ de plus en plus grand, on trouve un \uline{pro-motif} \inserted{(en alg de Lie)} canonique, appelé \uline{promotif fondamental}, qui peut aussi se construire de façon canonique en termes de la catégorie des motifs, munie de ses structures tensorielles \ldots

Quelle est la partie ``algébrique'' de ce promotif ? [Sur un corps fini, c'est tout -- le promotif n'étant lui-$=$ autre que $\mathbb Z(0)$ ! ; \deleted{dans}\inserted{en} tout \deleted{\unsure{autre}}\inserted{autre} cas, \deleted{\unsure{on}} \inserted{je pense qu'il est \uline{nul} \ldots.} \deleted{\unsure{il s'adresse|il s'intéresse} contenu dans le \uline{centre} du dit promotif fondamental, et je pense que c'est réduit à $\mathbb Z(0)$ [\unsure{n-} ``scalaires'']}
% archive sheet number: 58

\scanpage{60}
15) \uline{Cohomologie absolue}

En plus des foncteurs cohomologiques de motifs $\mathbb R^{\unsure{*|+}}f_*$ [$f$ morphisme de type fini] il y a lieu de considérer également des foncteurs ``absolus''
\[
H^*(X,M)=(H^i(X,M))_i
\]
$M$ motif variable sur $X$, qui doivent provenir d'un foncteur de catégories triangulées
\[
\mathbb R\Gamma_X:\ D^b(\deleted{\unsure{\mathcal M}}\mathcal M^{\deleted{\unsure{+}}}(X))\longrightarrow D^b(\mathrm{Ab})
\]
On veut que les $H^*(X,M)$ et $\mathbb R\Gamma(\unsure{M.|M_.|M_{\bullet}})$ aient les propriétés suivantes

$1^\circ)$ Commutation aux $\varprojlim$ essentiellement affines de préschémas . \unsure{qu.|q.} cpt \unsure{qu.|q.} sép. Cela nous ramène pratiquement aux cas $X$ de type fini sur $\mathbb Z$..

$2^\circ)$ Transitivité : \unsure{Si|si} $f:X\mathbin{-}Y$ de
% archive sheet number: 59

\scanpage{61}
% Faint isolated strokes in the middle, oriented as in the scan.
\handwritten{\figure{\begin{tikzpicture}[x=1cm,y=1cm]
\draw[black!25,very thin] (0,.24)--(-.025,.07);
\draw[black!55,thin,->] (1.25,0)--(.67,.02);
\draw[black!60,thin] (3.05,.12)--(3.01,-.07);
\draw[black!60,thin] (3.20,.12)--(3.21,.01)--(3.25,-.10);
\end{tikzpicture}}}

% Typescript rotated 180 degrees relative to the archive foliation.
\deleted{il en est de même de la famille $(u')$ qui s'en déduit par le changement de base $X'\quad X$.

On aura expliqué au préalable}
% archive sheet number: 60

\scanpage{62}
type fini, on a
\[
\mathbb R\Gamma_X=\mathbb R\Gamma_Y\,\mathbb R f
\]
d'où des suites spectrales \unsure{dans|ds} $(\mathrm{Ab})$
\[
H^*(X,M)\Longleftarrow H^p(Y,R\unsure{1|{}^q}f_*(M))
\]

$3^\circ)$ Bien entendu, on veut que $H^i(X,M)=0$ si $i<0$, donc $H^0(X,M)$ est exact à g. De plus, isom fonctoriel
\[
H^0(X,M)=\operatorname{Hom}(\deleted{\mathbb Z}\mathbb Q_X(0),M)
\]
[N.B. Rien n'empêche d'introduire des hyperext
\[
\mathbb E\mathrm{xt}^i(X;M,N)=\mathbb H^i(X,\mathbb R\mathcal H\mathrm{om}(M,N)).\ ]
\]
Cela redonne bien, par le n°~7
\[
\operatorname{Hom}(M,N)\simeq\mathbb E\mathrm{xt}^0(X,M,N)
\]
et on espère que $\mathrm{Ext}^{1}(X;M,N)$ est bien l'ensemble des \uline{classes} \uline{d'extension} de $M$ par $N$. [On aimerait pouvoir dire que
% archive sheet number: 61

\scanpage{63}
$\mathbb E\mathrm{xt}^*(X;M,N)$ est le \deleted{\unsure{gr|gr.}} $\mathrm{Ext}$ \unsure{dans|de} la catégorie des motifs, i.e. aux foncteurs dérivés ? \ldots.

Pourtant, je doute que les $\mathrm{Ext}^i$ $(i>0)$ soient effaçables\ldots. \inserted{[peut-être ds la catégorie des Ind-motifs ???]} \unsure{Noter|Notons} la suite spectrale
\[
\mathbb H^*(X,M^{\unsure{\cdot|\bullet}})\Longleftarrow H^p(X,H^q(M^\cdot))
\]
d'où en particulier
\[
\mathbb E\mathrm{xt}^*(X,M,N)\Longleftarrow \unsure{\mathbb H|H}^p(X,\underline{\mathrm{Ext}}^{\unsure{q|1}}(M,N))
\]

$4^\circ)$ On veut \deleted{que \unsure{sous des conditions standard} sur $X$} \inserted{des hom de $\delta$-foncteurs}
\[
H^*(X,M)\deleted{\otimes}_{\mathbb Q}\mathbb Q_\ell\longrightarrow H^*(X,T_\ell(M))
\]
d'où aussi
\[
\mathrm{Ext}^*(X;M,N)\otimes_{\mathbb Q}\mathbb Q_\ell\longrightarrow\mathbb E\mathrm{xt}^{\unsure{*|i}}(X;T_\ell(M),T_\ell(N))
\]
qui, \deleted{sous des conditions d'excellence sur $X$,} \inserted{sous l'hypothèse que $X$ est \uline{de type fini sur} $\mathbb Z$,} (\unsure{assurant|motivant|en invoquant} le fait que les seconds mbs sont de t. f. sur $\mathbb Q_\ell$)
% Three parallel emphasis strokes to the left of the last two lines.
\margin{\tikz[x=1cm,y=1cm]{\draw (0,0)--(-.03,.95); \draw (.10,.02)--(.06,1.05); \draw (.19,.04)--(.16,1.10);}}
% archive sheet number: 62

